\hypertarget{rearrangement-cycles-quantum-check-of-the-handoff-and-static-array-uniformity}{%
\chapter{Rearrangement cycles, quantum check of the handoff, and
static-array
uniformity}\label{rearrangement-rearrangement-cycles-quantum-check-of-the-handoff-and-static-array-uniformity}}

Follow-up to \texttt{docs/transport3d.md}. Code:
\texttt{picsim/rearrange.py}, \texttt{scripts/m20\_rearrange.py}
(scheduler), \texttt{scripts/m19\_quantum.py} (quantum check),
\texttt{scripts/m18\_transport3d.py} parts \texttt{uniformity} and
\texttt{slowlane}; results \texttt{results/m20\_rearrange.json},
\texttt{m19\_quantum.json}, \texttt{m18\_uniformity.json},
\texttt{m18\_slowlane.json}. Values are SIM and valid in the regime
stated. They should be reviewed before they are quoted.

\hypertarget{1-result}{%
\section{Result}\label{rearrangement-1-result}}

A 16-chip stack with row select assembles a defect-free 32 x 32 block
from a 50 x 45 array loaded at 50\% in 17 ms (mean of 40 loadings; worst
23 ms), using about 8 passes. Every loading succeeded before losses. The
average atom in the block moves 1.12 times; none moves more than twice.
With slower handoff steps (200 us ramps, 400 us lane steps), each move
heats the atom by 0.2 to 0.3 uK at a moving depth of only 0.3 mK. The
quantum and 3D classical models agree on this. The finished block is
therefore no hotter than about 0.6 uK above its starting temperature.
Loss sets the outcome: at 1e-3 loss per move, a cycle leaves 1.15
defects on average; at 1e-4, 0.12. Vacuum loss during the 17 ms adds
0.59 defects for a 30 s lifetime and 0.06 for 300 s. The static
array\textquotesingle s depth uniformity is not a constraint: 20\% rms
random depth errors cost nothing.

{\footnotesize
\begin{longtable}[]{@{}>{\raggedright\arraybackslash}p{0.091\dimexpr(\linewidth-14\tabcolsep)\relax}>{\raggedright\arraybackslash}p{0.086\dimexpr(\linewidth-14\tabcolsep)\relax}>{\raggedright\arraybackslash}p{0.171\dimexpr(\linewidth-14\tabcolsep)\relax}>{\raggedright\arraybackslash}p{0.126\dimexpr(\linewidth-14\tabcolsep)\relax}>{\raggedright\arraybackslash}p{0.149\dimexpr(\linewidth-14\tabcolsep)\relax}>{\raggedright\arraybackslash}p{0.189\dimexpr(\linewidth-14\tabcolsep)\relax}>{\raggedright\arraybackslash}p{0.189\dimexpr(\linewidth-14\tabcolsep)\relax}@{}}
\toprule\noalign{}
target (loading) & addressed rows & passes (vertical + horizontal) &
cycle time, mean / max & moves per atom, mean / max & defects from moves
at 1e-3 / 1e-4 & defects from vacuum at 30 / 300 s \\
\midrule\noalign{}
\endhead
\bottomrule\noalign{}
\endlastfoot
30 x 30 (50\%) & 16 & 4.3 + 2 & 13.3 / 17.1 ms & 1.11 / 2 & 1.00 / 0.10
& 0.40 / 0.04 \\
\textbf{32 x 32 (50\%)} & \textbf{16} & \textbf{5.8 + 2} & \textbf{17.2
/ 23.0 ms} & \textbf{1.12 / 2} & \textbf{1.15 / 0.12} & \textbf{0.59 /
0.06} \\
32 x 32 (60\%) & 16 & 5.0 + 2 & 14.8 / 17.1 ms & 1.03 / 2 & 1.06 / 0.11
& 0.50 / 0.05 \\
36 x 36 (60\%) & 16 & 8.3 + 3 & 25.1 / 35.6 ms & 1.07 / 2 & 1.38 / 0.14
& 1.08 / 0.11 \\
32 x 32 (50\%) & 32 (two stacks) & 5.8 + 1 & 15.5 / 21.3 ms & 1.12 / 2 &
1.15 / 0.12 & 0.53 / 0.05 \\
\end{longtable}
}

The probability that a 32 x 32 block is defect-free after one cycle is
0.84 at 1e-4 loss per move and a 300 s lifetime. Reaching 0.99 in one
cycle would need about 5e-6 loss per move and a vacuum lifetime above
about an hour; a repair round after imaging is the realistic route.

\hypertarget{2-scheduler}{%
\section{Scheduler}\label{rearrangement-2-scheduler}}

The hardware sets the rules: (i) a stack addresses 16 consecutive rows
at a time; (ii) every tweezer of the stack shares the y motion (into the
lane 2 um off the row, optionally along mid-column lanes by j rows, and
back), while each moves independently in x; (iii) tweezers in one row
keep their order and so stay at least 4 um apart.

\textbf{Vertical balancing.} Each pass chooses a stack position k and a
common row shift j that move the most surplus atoms into target rows
lacking atoms. Arriving atoms go straight to empty target columns where
possible. Atoms that arrive in a pass are not picked again in the same
pass.

\textbf{Horizontal compaction.} One pass per group of 16 target rows.
Each row\textquotesingle s atoms are matched in order to the target
columns, with extras pushed just outside the block.

A pass takes two ramps, two lane steps, the x moves (min-jerk; 400 us
per 40 um, scaled as sqrt(d), at least 200 us), the y travel for j rows,
and 0.3 ms of scanner settling (CHOICE). This greedy schedule is not
optimal. A planner that combines balancing and compaction, or shares
passes between stack positions, would need fewer passes.

\textbf{Sensitivity of the cycle time} (32 x 32, 50\%, one stack):

{\small
\begin{longtable}[]{@{}ll@{}}
\toprule\noalign{}
timing & cycle time \\
\midrule\noalign{}
\endhead
\bottomrule\noalign{}
\endlastfoot
base (ramps 100 us, lane steps 400 us) & 17.2 ms \\
low-heating (ramps 200 us, lane steps 400 us) & 18.8 ms \\
fast lane steps (100 us) & 12.5 ms \\
all fast (ramps 50, lanes 100, x 200 us per 40 um) & 8.2 ms \\
scanner settling 1 ms instead of 0.3 ms & 22.7 ms \\
\end{longtable}
}

Because vacuum loss grows with cycle time, the fast setting halves it
(0.28 defects at 30 s). It costs heating: about 10 uK per move at 0.3
mK, or 1.2 uK at 0.5 mK (Sec. 3). Hence two regimes: slow and cold at
0.3 mK (low power), or fast at 0.5 mK.

\hypertarget{3-quantum-check-of-the-handoff}{%
\section{Quantum check of the
handoff}\label{rearrangement-3-quantum-check-of-the-handoff}}

A 1D split-step propagation along y (the handoff and lane axis) starts
the atom in the ground state of the static site. The y profile is the
moving tweezer\textquotesingle s, taken from the 3D grid, with gravity
included. The x move is treated separately. Results are compared with
classical trajectories from a Wigner sample of the same ground state:

{\small
\begin{longtable}[]{@{}>{\raggedright\arraybackslash}p{0.148\dimexpr(\linewidth-10\tabcolsep)\relax}>{\raggedright\arraybackslash}p{0.182\dimexpr(\linewidth-10\tabcolsep)\relax}>{\raggedright\arraybackslash}p{0.227\dimexpr(\linewidth-10\tabcolsep)\relax}>{\raggedright\arraybackslash}p{0.193\dimexpr(\linewidth-10\tabcolsep)\relax}>{\raggedright\arraybackslash}p{0.250\dimexpr(\linewidth-10\tabcolsep)\relax}@{}}
\toprule\noalign{}
hold / moving & ramp / lane step & energy gain, quantum & ground state
kept & energy gain, classical \\
\midrule\noalign{}
\endhead
\bottomrule\noalign{}
\endlastfoot
0.1 / 0.3 mK & 100 / 100 us & 9.7 uK & 0.08 & 11.6 uK \\
0.1 / 0.3 mK & 100 / 200 us & 2.0 uK & 0.21 & 1.4 uK \\
0.1 / 0.3 mK & \textbf{200 / 400 us} & \textbf{0.12 uK} & \textbf{0.98}
& 0.94 uK \\
0.1 / 0.5 mK & 100 / 100 us & 0.84 uK & 0.84 & 1.53 uK \\
0.1 / 0.5 mK & 200 / 400 us & 0.13 uK & 0.99 & 0.87 uK \\
0.2 / 0.6 mK & 200 / 400 us & 0.15 uK & 0.99 & 1.72 uK \\
0.2 / 1.0 mK & 100 / 100 us & 1.58 uK & 0.66 & 1.58 uK \\
\end{longtable}
}

(Trap quantum of the static site: 1.4 uK at 0.1 mK, 2.0 uK at 0.2 mK.)
The x move adds 0.08 to 0.37 uK in both models (ground state kept 0.98
to 0.99). Two findings follow. (i) The lane step, not the x move, sets
the heating, and it falls steeply with its duration: from 100 to 400 us,
the gain at 0.1 / 0.3 mK drops 80-fold. (ii) Above about one trap
quantum the classical model agrees with the quantum one to within a
factor of 2, on the pessimistic side. Below it, classical Wigner
sampling overstates the heating by a factor of up to 10, as expected for
zero-point motion. A 3D classical run with the slow steps confirms the
delivery: 3,000 of 3,000 atoms, with 0.16 uK (5 uK atoms, 0.1 / 0.3 mK)
and 0.59 uK (10 uK, 0.2 / 0.6 mK) of heating, and no loss of waiting
atoms.

\hypertarget{4-static-array-depth-uniformity}{%
\section{Static-array depth
uniformity}\label{rearrangement-4-static-array-depth-uniformity}}

2,000 atoms per case, lane transport, three operating points
(results/m18\_uniformity.json):

{\small
\begin{longtable}[]{@{}>{\raggedright\arraybackslash}p{0.434\dimexpr(\linewidth-8\tabcolsep)\relax}>{\raggedright\arraybackslash}p{0.186\dimexpr(\linewidth-8\tabcolsep)\relax}>{\raggedright\arraybackslash}p{0.186\dimexpr(\linewidth-8\tabcolsep)\relax}>{\raggedright\arraybackslash}p{0.195\dimexpr(\linewidth-8\tabcolsep)\relax}@{}}
\toprule\noalign{}
non-uniformity & 5 uK, 0.1 / 0.3 mK & 5 uK, 0.1 / 0.5 mK & 10 uK, 0.2 /
0.6 mK \\
\midrule\noalign{}
\endhead
\bottomrule\noalign{}
\endlastfoot
source site deeper by 20 / 40 / 60\% & 1.000 / 1.000 / 0.996 & 1.000 /
1.000 / 1.000 & 1.000 / 1.000 / 1.000 \\
target site shallower, x 0.8 / 0.6 / 0.4 & 1.000 / 0.999 / 0.998 & 1.000
/ 1.000 / 1.000 & 1.000 / 1.000 / 0.9995 \\
random depth errors 5 / 10 / 20\% rms (two draws) & 1.000 all & 1.000
all & 1.000 all \\
\end{longtable}
}

The handoff tolerates the static array\textquotesingle s depth errors
far beyond what a meta-optic or SLM array would show. The energy figures
in \texttt{m18\_uniformity.json} include the adiabatic change between
unequal source and target depths and should not be read as heating.

\hypertarget{5-limits-and-next-steps}{%
\section{Limits and next steps}\label{rearrangement-5-limits-and-next-steps}}

The scheduler is greedy and uses one rearrangement cycle without repair.
The loss per move is not computed: the simulations bound it below about
1e-3 (0 of 3,000 atoms), and 1e-4 is an assumption. The quantum model is
1D per axis. Photon scattering by the trapping light (estimated at a few
per second at 0.3 to 0.5 mK) and imaging are not included; neither are
moves of atoms whose heating would make them leave the target sites
during later moves.

Next steps: (i) a repair round after imaging, with its own pass count
and time; (ii) an optimised schedule (passes shared across stack
positions); (iii) a 3D quantum or semiclassical check of the axial
motion during the handoff.

\hypertarget{verification-notes}{%
\section{Verification notes}\label{rearrangement-verification-notes}}

{\small
\begin{longtable}[]{@{}>{\raggedright\arraybackslash}p{0.462\dimexpr(\linewidth-8\tabcolsep)\relax}>{\raggedright\arraybackslash}p{0.092\dimexpr(\linewidth-8\tabcolsep)\relax}>{\raggedright\arraybackslash}p{0.069\dimexpr(\linewidth-8\tabcolsep)\relax}>{\raggedright\arraybackslash}p{0.377\dimexpr(\linewidth-8\tabcolsep)\relax}@{}}
\toprule\noalign{}
value & location & type & pointer \\
\midrule\noalign{}
\endhead
\bottomrule\noalign{}
\endlastfoot
passes, times, moves, defect estimates & Sec. 1, 2 & SIM &
results/m20\_rearrange.json (40 loadings; review) \\
quantum and Wigner energy gains & Sec. 3 & SIM &
results/m19\_quantum.json (review) \\
slow-lane 3D check & Sec. 3 & SIM & results/m18\_slowlane.json (3,000
atoms) \\
uniformity table & Sec. 4 & SIM & results/m18\_uniformity.json (2,000
atoms) \\
per-move loss 1e-3 / 3e-4 / 1e-4, vacuum lifetimes 30 / 300 s, scanner
settling 0.3 ms & Sec. 1, 2 & CHOICE & assumptions \\
\end{longtable}
}
